% !TSS name: Звіт з лабораторної роботи (укр.)
% !TSS description: звіт: мета, прилади, теорія, хід роботи, таблиця результатів (siunitx + booktabs), обробка похибок, висновки; LuaLaTeX, кирилиця (питає номер, назву, автора, групу)
% !TeX program = lualatex
% !TeX encoding = utf8
% !TeX spellcheck = uk_UA
\documentclass[a4paper, 12pt]{article}
\usepackage{fontsetup}
\usepackage[english, ukrainian]{babel}
\usepackage[a4paper, margin=2cm]{geometry}
\usepackage{amsmath}
\usepackage{graphicx}
\usepackage{booktabs}
\usepackage{siunitx}
\sisetup{output-decimal-marker={,}, separate-uncertainty}
% українські позначення одиниць (додай свої: \DeclareSIUnit{\kilogram}{кг})
\DeclareSIUnit{\meter}{м}
\DeclareSIUnit{\second}{с}
\usepackage{hyperref}

\title{Лабораторна робота № ${ask:num:1}\\ ${title}}
\author{${ask:author}, група ${ask:group}}
\date{${date}}

\begin{document}
\maketitle

\section{Мета роботи}
${cursor}

\section{Прилади й матеріали}
\begin{itemize}
	\item
\end{itemize}

\section{Теоретичні відомості}
Робоча формула:
\begin{equation}\label{eq:main}
	x = \ldots
\end{equation}

\section{Хід роботи}
\begin{enumerate}
	\item
\end{enumerate}

\section{Результати вимірювань}
\begin{table}[h!]
	\centering
	\caption{Результати вимірювань}\label{tab:data}
	\begin{tabular}{c S[table-format=2.2] S[table-format=1.3]}
		\toprule
		№ & {$x$, \si{\meter}} & {$t$, \si{\second}} \\
		\midrule
		1 & 10.00 & 1.000 \\
		2 & 20.00 & 2.000 \\
		\bottomrule
	\end{tabular}
\end{table}

\section{Обробка результатів і похибки}
Результат: $x = \SI{10.0 \pm 0.1}{\meter}$.

\section{Висновки}

\end{document}
